\documentclass[UTF8,11pt]{ctexart}
\usepackage[a4paper,margin=24mm]{geometry}
\usepackage{amsmath,amssymb,amsthm,mathtools,mathrsfs}
\usepackage[all]{xy}
\usepackage{xurl,hyperref,enumitem}
\usepackage{tikz}
\usetikzlibrary{arrows.meta,cd,decorations.markings,calc,patterns}
\hypersetup{hidelinks,breaklinks=true}
\setlength{\emergencystretch}{3em}
\allowdisplaybreaks
\newtheorem*{theorem}{Theorem}
\newtheorem*{thm}{Theorem}
\newtheorem*{lemma}{Lemma}
\newtheorem*{proposition}{Proposition}
\newtheorem*{prop}{Proposition}
\newtheorem*{corollary}{Corollary}
\newtheorem*{cor}{Corollary}
\newtheorem*{claim}{Claim}
\theoremstyle{definition}
\newtheorem*{definition}{Definition}
\newtheorem*{defn}{Definition}
\newtheorem*{remark}{Remark}
\newtheorem*{example}{Example}
\newcommand{\R}{\mathbb R}
\newcommand{\C}{\mathbb C}
\newcommand{\Z}{\mathbb Z}
\newcommand{\Q}{\mathbb Q}
\newcommand{\N}{\mathbb N}
\newcommand{\F}{\mathbb F}
\newcommand{\D}{\mathbb D}
\newcommand{\bB}{\mathbb B}
\newcommand{\bC}{\mathbb C}
\newcommand{\bR}{\mathbb R}
\newcommand{\bZ}{\mathbb Z}
\newcommand{\bN}{\mathbb N}
\newcommand{\bQ}{\mathbb Q}
\newcommand{\bD}{\mathbb D}
\newcommand{\bF}{\mathbb F}
\newcommand{\CP}{\mathbb{CP}}
\newcommand{\RP}{\mathbb{RP}}
\newcommand{\sO}{\mathscr O}
\newcommand{\sI}{\mathscr I}
\newcommand{\sF}{\mathscr F}
\newcommand{\sG}{\mathscr G}
\newcommand{\sH}{\mathscr H}
\newcommand{\sR}{\mathscr R}
\newcommand{\sM}{\mathscr M}
\newcommand{\sV}{\mathscr V}
\newcommand{\sU}{\mathscr U}
\DeclareMathOperator{\Hom}{Hom}
\DeclareMathOperator{\Ext}{Ext}
\DeclareMathOperator{\Tor}{Tor}
\DeclareMathOperator{\Spec}{Spec}
\DeclareMathOperator{\Proj}{Proj}
\DeclareMathOperator{\supp}{supp}
\DeclareMathOperator{\im}{im}
\DeclareMathOperator{\id}{id}
\DeclareMathOperator{\Tr}{Tr}
\DeclareMathOperator{\tr}{tr}
\DeclareMathOperator{\rank}{rank}
\DeclareMathOperator{\ord}{ord}
\DeclareMathOperator{\dist}{dist}
\DeclareMathOperator{\End}{End}
\DeclareMathOperator{\Aut}{Aut}
\DeclareMathOperator{\Gal}{Gal}
\DeclareMathOperator{\vol}{vol}
\DeclareMathOperator{\Cl}{Cl}
\DeclareMathOperator{\coker}{coker}
\DeclareMathOperator{\codim}{codim}
\DeclareMathOperator{\divergence}{div}
\DeclareMathOperator{\Ric}{Ric}
\DeclareMathOperator{\grad}{grad}
\DeclareMathOperator{\Hess}{Hess}
\newcommand{\abs}[1]{\left\lvert#1\right\rvert}
\newcommand{\sabs}[1]{\lvert#1\rvert}
\newcommand{\babs}[1]{\bigl\lvert#1\bigr\rvert}
\newcommand{\Babs}[1]{\Bigl\lvert#1\Bigr\rvert}
\newcommand{\bbabs}[1]{\biggl\lvert#1\biggr\rvert}
\newcommand{\BBabs}[1]{\Biggl\lvert#1\Biggr\rvert}
\newcommand{\norm}[1]{\left\lVert#1\right\rVert}
\newcommand{\snorm}[1]{\lVert#1\rVert}
\newcommand{\bnorm}[1]{\bigl\lVert#1\bigr\rVert}
\newcommand{\Bnorm}[1]{\Bigl\lVert#1\Bigr\rVert}
\newcommand{\bbnorm}[1]{\biggl\lVert#1\biggr\rVert}
\newcommand{\BBnorm}[1]{\Biggl\lVert#1\Biggr\rVert}
\newcommand{\widebar}[1]{\overline{#1}}
\newcommand{\nicefrac}[2]{\frac{#1}{#2}}
\newcommand{\linnprod}[2]{\langle#1,#2\rangle}
\newcommand{\myindex}[1]{#1}
\newcommand{\glsadd}[1]{}
\newcommand{\avoidbreak}{}
\newcommand{\sourceref}[1]{\textnormal{[原文：\nolinkurl{#1}]}}
\newcommand{\thmref}[1]{\sourceref{#1}}
\newcommand{\lemmaref}[1]{\sourceref{#1}}
\newcommand{\propref}[1]{\sourceref{#1}}
\newcommand{\corref}[1]{\sourceref{#1}}
\newcommand{\defnref}[1]{\sourceref{#1}}
\newcommand{\exampleref}[1]{\sourceref{#1}}
\newcommand{\exerciseref}[1]{\sourceref{#1}}
\newcommand{\figureref}[1]{\sourceref{#1}}
\newcommand{\sectionref}[1]{\sourceref{#1}}
\newcommand{\appendixref}[1]{\sourceref{#1}}
\newcommand{\stacksref}[2]{\href{https://stacks.math.columbia.edu/tag/#1}{#2 [#1]}}

\begin{document}
\section*{044\quad 四维分歧覆盖上的辛形式与同痕唯一性}
\subsection*{来源与计数目标}
Kartik Venkatram 与 Denis Auroux 合作整理的 MIT 18.966 \emph{Geometry of Manifolds} 人类课程讲义，Spring 2007。课程页面明确记载该署名关系；不是将学生笔记冒充授课者独著。Lecture 23，Definitions 1--3、局部分歧计算、Lemma 1、Proposition 1 及其完整证明，pp.~1--2。获取日期：2026-09-28。
\par\url{https://ocw.mit.edu/courses/18-966-geometry-of-manifolds-spring-2007/66f491c36656f58fae838a375df8e949_lect23.pdf}
\par\url{https://ocw.mit.edu/courses/18-966-geometry-of-manifolds-spring-2007/pages/lecture-notes/}
\par\textbf{本文件只计一个问题：}从含单分歧与尖点局部模型的辛分歧覆盖构造 $[\omega_X]=f^*[\omega_Y]$ 的辛形式，并证明本构造在同痕意义下的唯一性。
\subsection*{作者命题与人类证明}
\paragraph{1. Branched Covers}
\begin{definition}[1]
For $(M,\omega)$ a symplectic manifold, $p\in M$, a local diffeomorphism $\phi:U\to\C^n$ for $U\ni p$ is $\omega$-tame if $(\phi_*\omega)(v,iv)>0$ $\forall v\ne0\in\C^n$. This is, $\phi^*J_0$ is $\omega$-tame, i.e. complex lines in $\C^n$ give symplectic submanifolds in $M$.
\end{definition}
\begin{definition}[2]
A map $f:X^4\to(Y^4,\omega_Y)$ from a compact, oriented manifold to a compact, symplectic manifold is a symplectic branched covering if $\forall p\in X$, $\exists U\ni p$, $V\ni f(p)$ coordinate neighborhoods (with $\phi:U\to\C^2$ an oriented diffeomorphism, $\psi:V\to\C^2$ an $\omega_Y$-tame diffeomorphism) s.t. the right vertical map of the commutative diagram
\[
\begin{tikzcd}
X\supset U\arrow[r,"\phi"]\arrow[d,"f"']&\C^2\arrow[d,"\psi f\phi^{-1}"]\\
Y\supset V\arrow[r,"\psi"']&\C^2
\end{tikzcd}\tag{1}
\]
is one of
\begin{enumerate}
\item $(u,v)\mapsto(z_1,z_2)=(u,v)$ (local diffeomorphism),
\item $(u,v)\mapsto(z_1,z_2)=(u^2,v)$ (simple branching),
\item $(u,v)\mapsto(z_1,z_2)=(u^3-uv,v)$ (cusp).
\end{enumerate}
\end{definition}
\begin{definition}[3]
The ramification curve is the set $R\subset X$ s.t. $R=\{p\in X\mid df(p)\text{ not onto}\}$. The branch (discriminant) curve is $D=f(R)\subset Y$, i.e. $D=\{q\in Y\mid\#f^{-1}(q)<\deg f\}$.
\end{definition}
We can calculate these curves explicitly in local coordinates. For instance, in the case of simple branching, we have that $\operatorname{Jac}(f)=\det(df)=\left|\begin{smallmatrix}2u&0\\0&1\end{smallmatrix}\right|=2u$, so $R=\{u=0\}$, $D=\{z_1=0\}$. In the case of a cusp, we have
\[
\operatorname{Jac}(f)=\det(df)=\begin{vmatrix}3u^2-v&-u\\0&1\end{vmatrix}=3u^2-v\tag{2}
\]
so $R=\{v=3u^2\}$ and $D=\{27z_1^2=4z_2^3\}$. What happens at the cusp: $\forall p\in R$, $\ker df=\C\times\{0\}\subset T_p\C^2$, so $\ker df$ is transverse to $TR$ at most points, but not at the cusp.
\begin{lemma}[1]
$R\subset X$ is a smooth, 2-dimensional submanifold, and $D\subset Y$ is a symplectic submanifold of $Y$ immersed except at isolated points (corresponding to cusps). In local coordinates, $D$ is a complex curve, so $TD$ consists of complex lines and $\omega_Y|_{TD}>0$.
\end{lemma}
\begin{proposition}[1]
If $f:X^4\to(Y^4,\omega_Y)$ is a symplectic branched covering, then $X$ carries a symplectic form $\omega_X$ (canonical up to isotopy) s.t. $[\omega_X]=f^*[\omega_Y]$.
\end{proposition}
\begin{proof}
Note that $f^*\omega_Y$ is a closed 2-form which is nondegenerate outside of $R$. $\forall p\in R$, $K_p=\ker df_p\subset T_pX$ is the kernel of $f^*\omega_Y$, and is a complex line in local coordinates (so it carries a natural orientation). We claim that $\exists\alpha$ an exact 2-form on $X$ s.t. $\forall p\in R$, $\alpha|_{K_p}>0$ is positive nondegenerate. Assuming this, we have that $\omega_X=f^*\omega_Y+\epsilon\alpha$ for $\epsilon>0$ sufficient small is closed and nondegenerate, since
\[
\omega_X\wedge\omega_X=f^*\omega_Y\wedge f^*\omega_Y+2\epsilon f^*\omega\wedge\alpha+\epsilon^2\alpha\wedge\alpha\tag{3}
\]
with the first term $\geq0$ everywhere and nondegenerate outside $R$, the second term positive on $R$ (in local coordinates, $f^*\omega_Y=\frac{i\lambda}{2}(dv\wedge d\bar v)$ for some $\lambda>0$, and $\alpha|_{\C\times\{0\}}>0$), and the third term negligible for small $\epsilon$.

We are left to prove our claim. Fix $p\in R$, and choose local coordinates $(u,v)$ on $X$ of our model. Set $x=\operatorname{Re}(u)$, $y=\operatorname{Im}(u)$, and $\alpha_p=d(\chi_1(|u|)\chi_2(|v|)xdy)$, where $\chi_1,\chi_2$ are cutoff functions chosen s.t. $\forall q\in R\cap\operatorname{Supp}(\chi_2)$, $\chi_1\equiv1$. In local coordinates, we have that
\[
\forall(u,v)\in R\cap\operatorname{Supp}(\chi_2),\qquad
\alpha_p|_{K=\C\times\{0\}}=\chi_2(|v|)dx\wedge dy>0\tag{4}
\]
and $0$ outside $\operatorname{Supp}(\chi_2)$. Covering $R$ by small neighborhoods, and taking $\alpha$ to be the sum of these $\alpha_p$ gives the desired exact form.

Finally, to see that the choice of $\omega_X$ is canonical up to isotopy, note that the set of $\alpha$'s satisfying our claim (i.e. exact 2-forms s.t. $\alpha|_K>0$ along $R$) is convex. That is, we can find an $\epsilon$ sufficiently small s.t., for two such forms $\alpha_1,\alpha_2$,
\[
f^*\omega_Y+\epsilon\alpha_1,\quad f^*\omega_Y+\epsilon\alpha_2,\quad
f^*\omega_Y+\epsilon(t\alpha_1+(1-t)\alpha_2))\tag{5}
\]
are all symplectic.
\end{proof}
\subsection*{整理说明（非作者正文）}
两页原图均实际读取，交换图按原箭头重排。省去两段不参与证明的 Remark，以及没有给出完整证明的反向 Theorem 1。式(3) 中一次省略的下标 $Y$、式(5) 末尾多出的右括号按原页保留。$\operatorname{Supp}$ 与其内部的正性措辞照录；不把排印整理称为数学认证。

标准先修为微分形式的拉回与楔积、局部坐标、紧支集截断函数、四维二形式的楔平方非退化判别、紧致性下统一小参数选择。局部核方向的恰当二形式构造、全局求和、保持上同调类、沿分歧集修复退化及凸性同痕全在本文。难度依据是从映射的复局部奇点组织出全局辛结构，不是直接引用覆盖空间存在定理。
\subsection*{可修改的简短标注（非作者正文）}
$c$：分歧覆盖总空间的辛结构

$a$：复局部模型；Jacobian 与分歧曲线；拉回二形式的核；截断函数；恰当形式修正；楔平方正性；凸性同痕

\texttt{cross\_theory\_status = yes}。分歧映射的复线核提供一致正方向，局部恰当形式可在该方向补足退化；局部构造拼为全局同调类不变的辛形式，再由修正项凸性得到同痕。位置：Proposition 1，p.~2。

全部标注均为非穷尽、可修改初稿；未标注不表示未使用。
\subsection*{许可与修改记录}
材料来自 MIT OpenCourseWare，作者与课程见来源。按该站所示 CC BY-NC-SA 4.0 许可复用，整理部分沿用同一许可。\url{https://creativecommons.org/licenses/by-nc-sa/4.0/}。2026-09-28：助手转录题证、调整数学排版并加入来源、范围和初稿标注；不暗示作者或 MIT 认可本次整理。所留为本次题证转录摘录，不是原 PDF 下载件。
\end{document}
